% ==============================================================================
% Plantilla Institucional de Trabajo Terminal — UPIIZ IPN
% Archivo: ejemplos/ecuaciones.tex
% ==============================================================================
% Ejemplos prácticos para ecuaciones individuales, sistemas, matrices y unidades SI.
% ==============================================================================

\section{Guía de ejemplos para ecuaciones matemáticas}
\label{sec:ejemplo-ecuaciones}

% ------------------------------------------------------------------------------
% 1. ECUACIÓN INDIVIDUAL NUMERADA
% ------------------------------------------------------------------------------
Para una ecuación individual numerada y referenciada, utilice el entorno \texttt{equation}:

\begin{equation}
    F(s) = m \, s^2 X(s) + b \, s X(s) + k X(s)
    \label{eq:ejemplo-masa-resorte}
\end{equation}

\noindent La función de transferencia en lazo abierto a partir de la \eqref{eq:ejemplo-masa-resorte} se obtiene como $G(s) = \frac{X(s)}{F(s)} = \frac{1}{m s^2 + b s + k}$.

% ------------------------------------------------------------------------------
% 2. SISTEMA DE ECUACIONES ALINEADAS (align)
% ------------------------------------------------------------------------------
Para varias líneas alineadas con el signo igual:

\begin{align}
    \dot{x}_1(t) &= x_2(t) \label{eq:espacio-estados-1} \\
    \dot{x}_2(t) &= -\frac{k}{m} x_1(t) - \frac{b}{m} x_2(t) + \frac{1}{m} u(t) \label{eq:espacio-estados-2}
\end{align}

% ------------------------------------------------------------------------------
% 3. MATRICES Y VECTORES
% ------------------------------------------------------------------------------
Para la representación en variables de estado matricial:

\begin{equation}
    \begin{bmatrix}
        \dot{x}_1 \\
        \dot{x}_2
    \end{bmatrix}
    =
    \begin{bmatrix}
        0 & 1 \\
        -\frac{k}{m} & -\frac{b}{m}
    \end{bmatrix}
    \begin{bmatrix}
        x_1 \\
        x_2
    \end{bmatrix}
    +
    \begin{bmatrix}
        0 \\
        \frac{1}{m}
    \end{bmatrix}
    u(t)
    \label{eq:matriz-estados}
\end{equation}

% ------------------------------------------------------------------------------
% 4. UNIDADES FÍSICAS CON SIUNITX
% ------------------------------------------------------------------------------
Para escribir unidades correctamente según la norma internacional, emplee el comando \texttt{\textbackslash qty\{magnitud\}\{unidad\}}:
\begin{itemize}
    \item Resistencia y capacitancia: $R = \qty{4.7}{\kilo\ohm}$, $C = \qty{100}{\micro\farad}$.
    \item Frecuencia y voltaje: $f = \qty{50}{\kilo\hertz}$, $V = \qty{12.6}{\volt}$.
    \item Velocidad y aceleración: $v = \qty{1.85}{\meter\per\second}$, $a = \qty{9.81}{\meter\per\second\squared}$.
    \item Temperatura y ángulo: $T = \qty{45.2}{\degreeCelsius}$, $\theta = \ang{30}$.
\end{itemize}
